\chapter{Introduction}
\label{ch:introduction}

\section{Motivation}

Most thyroid nodules are benign, but deciding which ones need a biopsy is a daily clinical task.
Ultrasound risk systems guide this decision~\citep{tessler2017}, and fine-needle aspiration
settles it. Machine-learning models can estimate the risk of malignancy from an ultrasound image,
but a model that is right most of the time is of limited use if nobody knows which of its
decisions can be trusted.

This thesis studies a model that decides only where it can guarantee its error rate. Nodules it
calls benign must contain at most 5\% cancers, nodules it calls malignant at most 20\% benign ones,
and every other nodule is referred to a clinician. The guarantee is obtained with Learn then
Test~\citep{angelopoulos2021ltt}, a conformal risk-control method that needs no assumption on the
model, only calibration data from the same distribution as future cases.

The underlying classifier follows Ra et al.~\citep{ra2025}: radiomics features are extracted from
the nodule, written as text, and read by a fine-tuned language model. The approach is attractive
because a language model could use what feature names mean, and because its answers can be
queried repeatedly to measure how stable they are.

\section{Research questions}

\begin{enumerate}
  \item[RQ1] How well does a language model classify thyroid nodules from radiomics prompts, and
             how does it compare with a classical model on the same features?
  \item[RQ2] Which choices of the pipeline matter: the feature extraction, the way values are
             written, and the size of the model?
  \item[RQ3] Can Learn then Test certify automatic decisions at clinically meaningful error rates,
             and does agreement across perturbed queries help?
\end{enumerate}

\section{Contributions so far}

\begin{itemize}
  \item An adaptation of the radiomics-to-language-model pipeline of Ra et al.\ to thyroid
        ultrasound (TN3K), with a three-way triage layer calibrated by Learn then Test.
  \item A data audit of TN3K: the official calibration and test sources are not exchangeable, and
        the original feature extraction had three defects (few grey levels, size in pixels,
        multi-component masks). A pooled split and a revised extraction address them.
  \item Evidence that better features raise the language model's test AUC from 0.70 to 0.75, indications
        that the model uses the magnitude of the values, and that it stays about 0.05 below XGBoost on
        the same features.
  \item A negative result with its explanation: no automatic decision could be certified, because
        the model's errors are stable and its cleanest low-risk groups are too small for the
        required evidence. Perturbation votes do not beat the softmax probability.
\end{itemize}

\section{Structure of the thesis}

Chapter~\ref{ch:background} introduces the clinical and methodological background.
Chapter~\ref{ch:data} describes the data, its quality issues and the splits.
Chapter~\ref{ch:methods} presents the pipeline and the triage layer.
Chapter~\ref{ch:results} reports the experiments, organised by question, and
Chapter~\ref{ch:discussion} discusses them. Chapter~\ref{ch:future} lists the ongoing and planned
work.
